\documentclass[12pt]{article}
\usepackage[utf8]{inputenc}
\usepackage[T1]{fontenc}
\usepackage{hyperref}
\title{Taruma Language Resources for K-12 Curriculum Development}
\author{Compiled by Claude Sonnet 5 (claude-sonnet-5) for Andr\'as Kornai}
\date{}
\begin{document}
\maketitle
\section{Introduction}
Taruma (ISO 639-3 tdm) is a language isolate of south-western Guyana (Upper Takutu--Upper Essequibo), spoken by very few elderly people, reported as about three speakers; they are also fluent in Wapishana. The language has been reported extinct several times. Search results show academic documentation but no teaching materials.
\section{Primary Educational Resources}
No primers, textbooks or school materials were found in search results.
\section{Linguistic and Reference Materials}
\begin{itemize}
\item Wikipedia overview: \url{https://en.wikipedia.org/wiki/Taruma_language}
\item Leiden University student thesis with a first phonological description and a proposed orthography: \url{https://studenttheses.universiteitleiden.nl/access/item%3A3454579/view}
\item ``Tracing the history of the Taruma people through plants, their names and ...'' (Wageningen University): \url{https://research.wur.nl/en/publications/tracing-the-history-of-the-taruma-people-through-plants-their-nam/}
\item Appendices dataset for that study: \url{https://research.wur.nl/en/datasets/appendices-for-tracing-the-history-of-the-taruma-people-through-p/}
\item Naturalis repository item (found in search; content not verified): \url{https://repository.naturalis.nl/pub/800921/}
\item Verbix Endangered Languages page: \url{https://docs.verbix.com/EndangeredLanguages/Taruma}
\end{itemize}
\section{Religious and Community Resources}
None found.
\section{Digital and Audio Resources}
\begin{itemize}
\item Taruma Corpus (Endangered Languages Documentation Programme archive): \url{https://eldp.access.preservica.com/dk0708/}
\item Endangered Languages Project entry: \url{https://endangeredlanguages.com/lang/6969}
\end{itemize}
\section{Curriculum Development Context}
With about three elderly speakers, a K-12 curriculum would rest on revitalisation from the ELDP corpus and the phonological description. No school use is known, and the speakers' working language is Wapishana.
\end{document}